\documentclass[12pt, a4paper]{article}

% ---------- Packages ----------
\usepackage[utf8]{inputenc}
\usepackage[T1]{fontenc}
\usepackage[british]{babel}
\usepackage[margin=2.5cm]{geometry}
\usepackage{graphicx}
\usepackage{hyperref}
\hypersetup{
    colorlinks = true,
    linkcolor  = black,
    urlcolor   = blue,
    citecolor  = black
}

% ---------- Title page info ----------
\title{B201 Computer Science Lab \\ \large Portfolio Report}
\author{Sumaya Macamo GH1038497}
\date{\today}

\begin{document}

% ---------- Title page ----------
\maketitle
\thispagestyle{empty}
\newpage

% ---------- Table of contents ----------
\tableofcontents
\newpage

% ---------- Sections ----------

\section{Introduction}
The purpose of this project was to develop a professional Computer Science portfolio for a working student application, in which my technical skills were showcased as well as my ability to organize, document and present software projects.

The portfolio itself consists of a personal website hosted on GitHub Pages, a GitHub repository that contains the project files, a CV that was created using LaTeX and a Python programming exercise to showcase a technical skill that I developed during the course of the module.

\section{Links}
\begin{itemize}
    \item Portfolio Website: \url{https://sumaya-macamo.github.io/individual-project/}
    \item GitHub Repository: \url{https://github.com/Sumaya-Macamo/individual-project}
\end{itemize}

\section{Structure of the Website and Repository}
In the GitHub repository, the main webpage is contained in \texttt{index.html}, while the styling is stored inside the \texttt{css} folder. The \texttt{assets} folder contains all the images I used as well as other documents such as the CV PDF as well as the LaTeX source code. I decided to separate these files into different folders not only for easy access and to improve readability, but also because it makes any future updates to be easier.

The website in turn, is organized into several sections, including Home, About, Skills, Projects, Education and Contact, each with its own specific purpose. The homepage begins with a hero section that briefly introduces me and provides easy and immediate access to the most important parts of the website. The navigation bar at the top, also allows people to move in between sections quickly without having to manually scroll down the entire website to find each section.

\section{Design Decisions}
My main goal with the portfolio was to create one that is both professional looking and also reflected a bit of who I am as person, or my personality in this case. I specifically chose the pastel pink colour for two main reasons: first it is one of my favourite colours, and second because I was inspired by colourful flowers, specifically tulips and roses, which I associate with creativity, growth and the ability to see beauty even in the simplest of things. I have a strong eye for aesthetics and visuals, and I try to evolve and learn every day, so I believed that the flower idea would represent that really well. I also used the same colours in my CV to create a consistent visual identity across the entire project.

The website uses a clean single-page layout with simple navigations so that it is easy to find information quickly without having to open multiple pages. I also chose project cards instead of a plain list in order to distinguish and highlight each project individually. The Education section was intentionally kept brief because the CV already contains the full details, which avoids unnecessary repletion. Finally, I developed the website using only HTML and CSS instead of a framework so that I could increase and strengthen my understating of the fundamentals, as well as have complete control over the design.

\section{Tools and Technologies}
To structure the website I used HTML, while CSS was used for styling, layout and its overall responsiveness. To host the project's source code and to track the changes made during its development I used Git and GitHub, while GitHub Pages was used to publish the website online. My CV was created using a LaTeX template in Overleaf. Finally, I also used Python to complete a small programming exercise that demonstrate my basic programming concepts.

\section{Technical Exercise}
As previously mentioned, the portfolio also includes a Python exercise. The Python program I developed essentially generates a lighter version of a hexadecimal colour by converting it to RGB values and then adjusting its brightness. I chose this exercise specifically because I wanted to challenge myself in regards to my python skills, and because it connects with the colour palette I used throughout the entire project. Although it is a relatively simple program, it demonstrates my understanding of functions, calculation and code organization.

\section{Reflection}
One of the main strengths of this portfolio is its clear organization and cohesiveness. All the information is easy to find, and the CV and website share the same visual identity, which creates a consistent personal brand. The project also demonstrates different sets of technical skills, such as web development, LaTeX and Python.

However, there are still some areas for improvement. First, the website currently has very little interactivity because I only used HTML and CSS. In the future as my skills develop, I would like to add JavaScript to improve the user experience. Additionally, I would also like to expand on the Projects section as I complete more coursework and develop personal projects.

\section{Conclusion}
Overall, developing this portfolio allowed me to combine several Computer Science skills that I developed into one project. It also demonstrates my current knowledge of web development, documentation and programming in general, and also provides a solid foundation where I can continue to improve as I gain more experience.

% ---------- References ----------
% Add any additional sources you used in Harvard style below.
\section*{References}
\begin{itemize}
    \item Kostyrka, A.V. (n.d.) \textit{LuxSleek-CV}. Overleaf. Available at: \url{https://www.overleaf.com/latex/templates/luxsleek-cv/qbvbqmrzxwyj} (Accessed: 26 June 2026).
\end{itemize}

\end{document}